\chapter{Introduction}
\label{ch:introduction}

Representing three-dimensional shapes in computers is a fundamental problem. This challenge appears everywhere: in medical imaging, where doctors need accurate models of organs for surgical planning; in robotics, where machines must understand the objects around them; and in entertainment, where realistic characters and environments bring virtual worlds to life. As these applications become more demanding, the need for better shape representation methods grows.

Traditional methods store shapes using discrete structures. Voxel grids, for example, divide space into small cubes and record whether each cube is occupied. This approach is simple to understand but expensive in practice. Doubling the resolution requires eight times more storage. A reasonably detailed shape at 512 resolution needs over 134 million data points. Point clouds take a different approach by storing only surface locations, which saves memory but loses information about connectivity. Polygon meshes represent surfaces as networks of triangles and work well for many applications, but they struggle when shapes have complex topology or need different levels of detail in different regions.

In recent years, researchers have developed a new approach using neural networks. Instead of storing discrete samples, a neural network learns a continuous function that describes the shape. Given any point in space, the network outputs whether that point is inside or outside the object. This idea, introduced by Mescheder et al.~\cite{mescheder2019occupancy}, offers two major advantages. First, the representation is resolution-independent: we can query the function at any precision without retraining. Second, storage depends only on network size, not output resolution. A network with a few hundred thousand parameters can represent shapes that would need billions of voxels.

However, standard neural networks have a significant limitation. They struggle to capture sharp edges and fine details. When trained to represent a complex shape, they produce smooth, rounded outputs that miss important geometric features. Rahaman et al.~\cite{rahaman2019spectral} studied this problem and called it spectral bias. Neural networks naturally learn smooth, low-frequency patterns first. High-frequency details, like sharp edges and fine textures, are much harder for them to capture. The mathematical explanation involves how gradients flow during training~\cite{jacot2018ntk}. In simple terms, the network's learning process favors gradual changes over rapid ones.

This limitation motivated researchers to develop new network architectures. Sitzmann et al.~\cite{sitzmann2020siren} proposed using sine functions as activations instead of the standard choices. Their method, called SIREN, naturally represents oscillating patterns and captures fine details much better than conventional networks. Saragadam et al.~\cite{saragadam2023wire} introduced WIRE, which uses Gabor wavelets. These mathematical functions combine oscillation with localization, allowing the network to focus on specific regions and frequencies. Ramasinghe and Lucey~\cite{ramasinghe2022gauss} experimented with Gaussian activations, which provide good spatial localization. Liu et al.~\cite{liu2024finer} developed FINER, where the activation frequency adapts based on the input. This allows the network to use different frequencies in different regions as needed.

Other researchers took different approaches. Tancik et al.~\cite{tancik2020fourier} showed that transforming input coordinates through random Fourier features helps standard networks learn high-frequency patterns. Kazerouni et al.~\cite{kazerouni2024incode} proposed INCODE, which uses a separate small network to adjust the main network's behavior during training. Fathony et al.~\cite{fathony2021mfn} introduced multiplicative filter networks, which combine layer outputs through multiplication rather than addition. Shi et al.~\cite{shi2024fr} developed a method that encodes frequency information directly into the network weights.

With so many methods available, choosing the right one becomes difficult. Essakine et al.~\cite{essakine2024inr} addressed this by conducting a comprehensive survey. They tested these methods on various tasks and identified three properties that seem to matter for performance. Spatial compactness refers to how well a method represents local features without affecting distant regions. Frequency compactness describes the ability to capture specific frequency ranges precisely. Adaptivity means the network can adjust its behavior based on what it needs to represent. This framework helps us understand the differences between methods, but important questions remain.

For three-dimensional shape reconstruction specifically, we still lack clear guidance. Which of these properties matters most when the goal is to capture sharp object boundaries? How much do the different methods cost in terms of computation time? When does a method fail, and can we predict this from its design? These practical questions are important for anyone trying to use these techniques in real applications. Without clear answers, practitioners must rely on trial and error, which wastes time and computational resources.

This thesis aims to answer these questions. We systematically evaluate eight neural network architectures for three-dimensional occupancy reconstruction. Our goal is to understand which design choices lead to better results and why. We build on the framework from Essakine et al. and test whether their identified properties explain performance differences in our specific task.

We focus on four main questions. First, how do the eight methods compare in reconstruction quality? Second, which architectural properties best predict good performance? Third, what causes methods to fail, and what do these failures look like? Fourth, how do training time and quality trade off against each other?

This work makes four contributions. First, we implement all eight methods for three-dimensional occupancy prediction under a single controlled protocol, run each over four random seeds, and make the code, the exact dataset, and the scripts that generate our tables and figures publicly available. Second, using error bars and a significance test, we identify a genuine best tier for this task---Fourier Features and INCODE, statistically tied with each other and significantly ahead of the rest---and show that a single run had instead suggested a different leader and a spurious near-tie, making the case concretely for repeated runs. Third, we use the gap between training and evaluation accuracy to separate three distinct failure modes---memorization, underfitting, and encoder misconfiguration---and show that the two weakest methods fail to fit even their training data, which points to optimization difficulty rather than the representational limits usually invoked to explain such results. Fourth, we show that architecture is not the largest lever: a single frequency-scale hyperparameter within one method outweighs the entire gap between architectures, and the survey properties we test do not predict the winners.

The rest of this thesis is organized as follows. Chapter 2 covers background material, including a detailed explanation of spectral bias and mathematical descriptions of each architecture. Chapter 3 describes our experimental setup: how we prepared data, trained networks, and measured results. Chapter 4 explains the implementation details and adaptations needed for three-dimensional inputs. Chapter 5 presents our experimental results with tables, comparisons, and analysis. Chapter 6 interprets those results, states the limitations that bound them, and concludes.
